% Lösungen Einheit 3 – Nr. 32–44
\einheitenkopf[e3]{Einheit 3 von 4 · Normalform und p-q-Formel}

\erg{32}{a) Wurzel ziehen (kein $x$-Glied) \quad b) Nullprodukt \quad c) rückwärts rechnen (quadrierte Klammer) \quad d) p-q-Formel \quad e) p-q-Formel (Normalform)}
\erg{33}{a) $p = 9$, $q = 4$ \quad b) $p = -5$, $q = 2$ \quad c) $p = 2$, $q = -8$ \quad d) $p = -1$, $q = -20$ \quad e) $p = -7$, $q = 0$}
\erg{34}{a) nein \quad b) ja, durch 2 \quad c) ja, durch $-1$ \quad d) ja, durch 5 \quad e) nein}
\erg{35}{a) $x_{1,2} = -3 \pm\sqrt{9 - 8} = -3 \pm 1$; $x_1 = -2$, $x_2 = -4$ \quad b) $-5 \pm 2$; $x_1 = -3$, $x_2 = -7$ \quad c) $-6 \pm 4$; $x_1 = -2$, $x_2 = -10$ \quad d) $-3 \pm 2$; $x_1 = -1$, $x_2 = -5$ \quad e) $1 \pm 3$; $x_1 = 4$, $x_2 = -2$}
\erg{36}{a) $x^2 - 6x + 5 = 0$; $x_1 = 5$, $x_2 = 1$ \quad b) $x^2 - 2x - 3 = 0$; $x_1 = 3$, $x_2 = -1$ \quad c) $x^2 + 2x - 8 = 0$; $x_1 = 2$, $x_2 = -4$ \quad d) $x^2 - 6x + 8 = 0$; $x_1 = 4$, $x_2 = 2$ \quad e) $x^2 + x - 6 = 0$; $-0{,}5 \pm\sqrt{6{,}25}$; $x_1 = 2$, $x_2 = -3$}
\erg{37}{a) $3 \pm\sqrt{0}$; $x = 3$ (eine Lösung) \quad b) $1 \pm\sqrt{-4}$; keine Lösung \quad c) Diskriminante $2{,}25 - 2 = 0{,}25$; zwei Lösungen \quad d) $6{,}25 - 7 = -0{,}75$; keine Lösung \quad e) $12{,}25 - 12{,}25 = 0$; eine Lösung}
\erg{38}{a) $x_{1,2} = 2 \pm\sqrt{3}$; $x_1 \approx 3{,}73$, $x_2 \approx 0{,}27$ \quad b) $x_{1,2} = -0{,}5 \pm\sqrt{1{,}25}$; $x_1 \approx 0{,}62$, $x_2 \approx -1{,}62$ \quad c) $x^2 - 2x - 1 = 0$; $x_{1,2} = 1 \pm\sqrt{2}$; $x_1 \approx 2{,}41$, $x_2 \approx -0{,}41$}
\erg{39}{a) $x_1 = 6$, $x_2 = -4$; Probe mit 6: $36 - 12 - 24 = 0$ (wA) \quad b) $x^2 - 4x - 5 = 0$; $x_1 = 5$, $x_2 = -1$; Probe in der Ausgangsgleichung mit $-1$: $3 + 12 - 15 = 0$ (wA)}
\erg{40}{a) $x^2 + 3x - 18 = 0$; $x_1 = 3$, $x_2 = -6$ \quad b) $x^2 + 2x + 1 = 3x + 7$, $x^2 - x - 6 = 0$; $x_1 = 3$, $x_2 = -2$ \quad c) $x^2 + 2x - 8 = 3x + 4$, $x^2 - x - 12 = 0$; $x_1 = 4$, $x_2 = -3$}
\erg{41}{a) Wurzel ziehen: $x^2 = 25$; $x_1 = 5$, $x_2 = -5$ \quad b) Ausklammern: $x_1 = 0$, $x_2 = -7$ \quad c) rückwärts: $x - 4 = \pm 1$; $x_1 = 5$, $x_2 = 3$ \quad d) Formel: $x_1 = 4$, $x_2 = -1$ \quad e) $-x^2 + 4x + 12 = 0$, $x^2 - 4x - 12 = 0$; $2 \pm\sqrt{16}$; $x_1 = 6$, $x_2 = -2$}
\erg{42}{a) Nicht normiert. Erst durch 3 teilen: $x^2 + 2x - 3 = 0$; $-1 \pm 2$; $x_1 = 1$, $x_2 = -3$. \quad b) Vorzeichen von $-\frac{p}{2}$: $p = -8$, also $-\frac{p}{2} = 4$; $4 \pm 2$; $x_1 = 6$, $x_2 = 2$. \quad c) Unter der Wurzel steht $(\frac{p}{2})^2 - q$, also $25 - 16 = 9$; $5 \pm 3$; $x_1 = 8$, $x_2 = 2$.}
\erg{43}{a) Die Wurzel aus einer negativen Zahl gibt es nicht, weil kein Quadrat negativ ist. \quad b) Die p-q-Formel gilt nur für die Normalform mit der Vorzahl 1 vor $x^2$; ohne Teilen kommen falsche Lösungen heraus (richtig: $x^2 - 5x + 6 = 0$; $x_1 = 3$, $x_2 = 2$).}
\erg{44}{a) $x^2 + 2x - 3 = 0$; $x_1 = 1$, $x_2 = -3$; $S_1(1|{-1})$, $S_2(-3|7)$ \quad b) am Graphen dieselben Punkte \quad c) $x^2 - 2x - 3 = 3x + 3$, $x^2 - 5x - 6 = 0$; $x_1 = 6$, $x_2 = -1$; $S_1(6|21)$, $S_2(-1|0)$}
